\subsection{Limitations}
\label{Subsection:Limitations}

The most important limitation is stated first because it bounds everything else. The performance figures reported in Section~\ref{Section:ExperimentalResults} are measured over a window that includes the years the agents trained on. The procedure does hold a year back from training and from the choice among an agent's own episodes, but the choice of which trained agent to publish was made on the full window, so the headline returns are an in-sample result. On the one year that no part of the procedure saw, the models lost money: the portfolio returned $-9.8\%$ in 2025 and five of the seven pairs were negative. That year contained a broad dollar reversal the training decade did not, and the one model that held up is the least directional of the set, but a single year is a small sample and nothing stronger should be read into it in either direction.

Selection is uncorrected in a second sense. Each published model was chosen from between $128$ and $256$ candidates, so a permutation $p$-value near $0.05$ is close to what that many draws produce by chance. The screening is therefore reported as a filter that demonstrably rejected bad models, not as a significance claim for the ones that survived. Under a strict correction only USD/CAD and GBP/USD would clearly stand.

Three narrower limitations belong on the record. Position size was chosen per pair after results were seen; it scales every trade without changing the policy that generates them, but it is a parameter fitted on the same data as everything else. Training is not reproducible across environments: it is bit-exact within one installation at a single thread, and a change of numerical library version alters the outcome, so the delivered models are archived artifacts with a verified replay rather than a recipe that regenerates them. And two weeks of January 2016 are missing from all seven pairs for reasons on the broker's side; the gap is uniform across pairs and falls outside the evaluation window, so it does not bias any comparison made here.

Finally, the scope is narrower than the method. Seven instruments in one asset class over one eleven-year period, with one decision cadence and one cost model, is a single point in a large space. Nothing in this work establishes that the design transfers to other markets, other periods or other holding horizons, and the USD/JPY result shows that it does not transfer to every instrument even within the seven.